% Generated draft: requires mathematical and source review.
\documentclass[11pt]{article}
\usepackage[margin=25mm]{geometry}
\usepackage{fontspec}
\setmainfont{Latin Modern Roman}
\usepackage{amsmath,amssymb,mathtools,unicode-math}
\setmathfont{Latin Modern Math}
\usepackage{enumitem,needspace,xurl,hyperref}
\setlength{\parindent}{0pt}
\setlength{\parskip}{5pt}
\setlength{\emergencystretch}{3em}
\newcommand{\N}{\mathbb{N}}
\newcommand{\Z}{\mathbb{Z}}
\newcommand{\Q}{\mathbb{Q}}
\newcommand{\R}{\mathbb{R}}
\newcommand{\C}{\mathbb{C}}
\newcommand{\F}{\mathbb{F}}
\newcommand{\A}{\mathbb{A}}
\newcommand{\PP}{\mathbb{P}}
\newcommand{\E}{\mathbb{E}}
\newcommand{\eps}{\varepsilon}
\newcommand{\dd}{\,\mathrm{d}}
\begin{document}
\section*{UnsolvedMath ID: 30000689}
\textbf{Source problem code:} catalogue entry
\subsection{Definitions and mathematical statement}
Let $d\geq1$ be an integer and let $\Sigma$ be a finite triangulation of $S^{2d}$. Its $d$-skeleton is
\[
 K=\Sigma^{(d)}=\{\sigma\in\Sigma:\dim\sigma\leq d\}.
\]
A missing $d$-face is a $(d+1)$-element subset $T\subseteq V(K)$ which is not a simplex of $K$ but whose every proper face is in $K$. Equivalently the boundary complex $\partial T$ is contained in $K$. Put
\[
 L=K\cup\{T\}=K\cup2^T,
\]
where the equality uses the existing boundary; the added face includes its already present subfaces and no new vertices.

The source conjectures
\[
 \forall\Sigma\ \forall T\text{ as above}:\qquad
 |L|\text{ does not embed in }S^{2d}.
\]
In the topological interpretation, an embedding means a continuous injective map, hence a homeomorphism onto its image for this compact complex. The conclusion concerns \emph{any} embedding of the augmented complex, allowing a new embedding of $K$; it is stronger than merely saying that the original skeleton embedding cannot be extended across the new simplex.

The source does not specify topological versus PL embedding or whether ``triangulated sphere'' is required to be a PL/combinatorial sphere. Those categories should be fixed before applying an obstruction. No polytopality or straight-line requirement is added. The positive-$d$ range is the usual nondegenerate setting; at $d=0$ a missing vertex on the existing vertex set does not occur.

The literature triage checked 21 August 2026 reports no authoritative complete resolution of the extracted claim. It supplies only the report and no independent obstruction theorem or case list; its dated open label is retained without current-status certification.
\subsection{Short English statement}
Take the middle-dimensional skeleton of a triangulated even-dimensional sphere and add a missing simplex whose boundary is already present. Must the augmented complex fail to embed in the same sphere, even after moving the original skeleton? The source leaves the topological/PL category to be confirmed.
\subsection{Sources}
\begin{itemize}
\item Local UnsolvedMath registry, permanent ID 30000689, exact source code catalogue entry; \emph{Embedding Obstructions from Missing Simplicial Faces}. Proposed year 2007, status open; corrected\_verified extraction retains the missing-face hypothesis and full nonembedding claim.
\item \emph{Geometric and Topological Combinatorics (2007)}, Oberwolfach Reports, DOI 10.4171/owr/2007/04, pp. 195--272, \url{https://doi.org/10.4171/owr/2007/04}.
\item Supplied literature triage checked 2026-08-21 cites the same report, \url{https://doi.org/10.4171/owr/2007/04}, and reports no authoritative complete resolution found. The dated OWR-derived curation requests independent verification of time-sensitive status and cited claims; no separate theorem/reference is supplied.
\end{itemize}
\end{document}
